\documentclass[10pt,a4paper]{amsart}
\usepackage{amsmath,amssymb,amsthm,graphicx}
\usepackage{verbatim}
\newenvironment{reference}{\comment}{\endcomment}
%\newenvironment{reference}{}{}
\newenvironment{slogan}{\comment}{\endcomment}
\newenvironment{history}{\comment}{\endcomment}
% For commutative diagrams we use Xy-pic
\usepackage[all]{xy}

% We use 2cell for 2-commutative diagrams.
\xyoption{2cell}
% Theorem environments.
%
\theoremstyle{plain}
\newtheorem{theorem}[subsection]{Theorem}
\newtheorem{proposition}[subsection]{Proposition}
\newtheorem{lemma}[subsection]{Lemma}

\theoremstyle{definition}
\newtheorem{definition}[subsection]{Definition}
\newtheorem{example}[subsection]{Example}
\newtheorem{exercise}[subsection]{Exercise}
\newtheorem{situation}[subsection]{Situation}

\theoremstyle{remark}
\newtheorem{remark}[subsection]{Remark}
\newtheorem{remarks}[subsection]{Remarks}

\numberwithin{equation}{subsection}

% Macros
%
\def\lim{\mathop{\mathrm{lim}}\nolimits}
\def\colim{\mathop{\mathrm{colim}}\nolimits}
\def\Spec{\mathop{\mathrm{Spec}}}
\def\Hom{\mathop{\mathrm{Hom}}\nolimits}
\def\Ext{\mathop{\mathrm{Ext}}\nolimits}
\def\SheafHom{\mathop{\mathcal{H}\!\mathit{om}}\nolimits}
\def\SheafExt{\mathop{\mathcal{E}\!\mathit{xt}}\nolimits}
\def\Sch{\mathit{Sch}}
\def\Mor{\mathop{\mathrm{Mor}}\nolimits}
\def\Ob{\mathop{\mathrm{Ob}}\nolimits}
\def\Sh{\mathop{\mathit{Sh}}\nolimits}
\def\NL{\mathop{N\!L}\nolimits}
\def\CH{\mathop{\mathrm{CH}}\nolimits}
\def\proetale{{pro\text{-}\acute{e}tale}}
\def\etale{{\acute{e}tale}}
\def\QCoh{\mathit{QCoh}}
\def\Ker{\mathop{\mathrm{Ker}}}
\def\Im{\mathop{\mathrm{Im}}}
\def\Coker{\mathop{\mathrm{Coker}}}
\def\Coim{\mathop{\mathrm{Coim}}}

% Boxtimes
%
\DeclareMathSymbol{\boxtimes}{\mathbin}{AMSa}{"02}

%
% Macros for moduli stacks/spaces
%
\def\QCohstack{\mathcal{QC}\!\mathit{oh}}
\def\Cohstack{\mathcal{C}\!\mathit{oh}}
\def\Spacesstack{\mathcal{S}\!\mathit{paces}}
\def\Quotfunctor{\mathrm{Quot}}
\def\Hilbfunctor{\mathrm{Hilb}}
\def\Curvesstack{\mathcal{C}\!\mathit{urves}}
\def\Polarizedstack{\mathcal{P}\!\mathit{olarized}}
\def\Complexesstack{\mathcal{C}\!\mathit{omplexes}}
% \Pic is the operator that assigns to X its picard group, usage \Pic(X)
% \Picardstack_{X/B} denotes the Picard stack of X over B
% \Picardfunctor_{X/B} denotes the Picard functor of X over B
\def\Pic{\mathop{\mathrm{Pic}}\nolimits}
\def\Picardstack{\mathcal{P}\!\mathit{ic}}
\def\Picardfunctor{\mathrm{Pic}}
\def\Deformationcategory{\mathcal{D}\!\mathit{ef}}
\usepackage{cleveref}
\crefname{originalreference}{original reference}{original references}
\Crefname{originalreference}{Original reference}{Original references}
\makeatletter
\newlabel{algebra-lemma-affine-blowup-quotient-description}{{algebra-lemma-affine-blowup-quotient-description}{source}{}{sourcepointer-1}{}}
\newlabel{algebra-lemma-affine-blowup-quotient-description@cref}{{[originalreference][0][]algebra-lemma-affine-blowup-quotient-description}{[1][0][]source}{}{}{}}
\makeatother
\crefname{originalreference}{original reference}{original references}
\Crefname{originalreference}{Original reference}{Original references}
\makeatletter
\newlabel{derived-lemma-trick-vanishing-composition}{{derived-lemma-trick-vanishing-composition}{source}{}{sourcepointer-2}{}}
\newlabel{derived-lemma-trick-vanishing-composition@cref}{{[originalreference][0][]derived-lemma-trick-vanishing-composition}{[1][0][]source}{}{}{}}
\makeatother
\title{Tor modules of an affine blowup algebra are annihilated by a power of the ideal in degrees beyond the number of generators}
\author{The Stacks Project Authors}
\date{stacks-project:local-cohomology.tex:local-snapshot-20261012; GNU FDL 1.2 or later}
\begin{document}
\maketitle
\noindent\textit{Editorial note.} All delivered mathematical text is the author's own, copied from the pinned Stacks Project chapter source sources/stacks-local-cohomology/pinned-original-source.tex; no statement or proof step was rewritten, reordered or removed.\par\medskip
\noindent\textit{Editorial note.} Preamble provenance: the Stacks Project keeps one single shared preamble for every chapter; this card takes it directly from the staged sources/stacks-local-cohomology/preamble.tex as three verbatim spans (lines 8-12, 14-18 and 154-217), so no derived file was written and nothing was added or redefined. The wrapper class and the base packages are supplied by the assembler.\par\medskip
\noindent\textit{Editorial note.} Bibliography provenance: neither the counted statement nor its proof contains a citation (this lane rescans the whole span before writing the card), so the card carries no bibliography and no entry was reconstructed.\par\medskip
\noindent\textit{Editorial note.} Selection provenance: this card was opened under the file-based control contract (CONTROL.md), with the mechanical pre-card gate run on its statement and proof interval before the card was written.\par\medskip
\noindent\textit{Allowed standard prerequisites (editorial).}\begin{itemize}
\item The affine blowup algebra of an ideal with respect to one of its elements, described as a quotient by the affine blowup relations modulo torsion. Invoked inside the counted proof at source lines 4562-4566 and attributed to a lemma of the cited chapter (a cross-chapter pointer). Provides the presentation of the algebra whose Tor modules the statement bounds.
\item The Koszul complex of the blowup relations and the identification of Tor with the homology of the derived tensor product. Used inside the counted proof at source lines 4566-4602. Reduces the annihilation statement to a statement about the homology of an explicit complex in high degrees.
\end{itemize}\medskip
\ifdefined\hypertarget\hypertarget{sourcepointer-1}{}\fi
\noindent\textit{Source reference.} Original reference algebra-lemma-affine-blowup-quotient-description: sources/stacks-local-cohomology/pinned-original-source.tex, source lines 4566--4566. Navigation pointer only: the counted text refers to the Stacks item algebra-lemma-affine-blowup-quotient-description, which belongs to another chapter of the project and is not part of this card. The pointer carries the target's own label as its number because this chapter's staged snapshot does not contain the other chapter's tag table; its content is not reproduced here.\par\medskip
\ifdefined\hypertarget\hypertarget{sourcepointer-2}{}\fi
\noindent\textit{Source reference.} Original reference derived-lemma-trick-vanishing-composition: sources/stacks-local-cohomology/pinned-original-source.tex, source lines 3675--3675. Navigation pointer only: the counted text refers to the Stacks item derived-lemma-trick-vanishing-composition, which belongs to another chapter of the project and is not part of this card. The pointer carries the target's own label as its number because this chapter's staged snapshot does not contain the other chapter's tag table; its content is not reproduced here.\par\medskip
\begin{lemma}
\label{lemma-annihilates-affine}
Let $I = (a_1, \ldots, a_t)$ be an ideal of a Noetherian ring $A$.
Set $a = a_1$ and denote $B = A[\frac{I}{a}]$ the affine blowup algebra.
There exists a $c > 0$ such that $\text{Tor}_i^A(B, M)$ is annihilated
by $I^c$ for all $A$-modules $M$ and $i \geq t$.
\end{lemma}

\begin{proof}
Recall that $B$ is the quotient of
$A[x_2, \ldots, x_t]/(a_1x_2 - a_2, \ldots, a_1x_t - a_t)$
by its $a_1$-torsion, see
Algebra, Lemma \ref{algebra-lemma-affine-blowup-quotient-description}. Let
$$
B_\bullet = \text{Koszul complex on }a_1x_2 - a_2, \ldots, a_1x_t - a_t
\text{ over }A[x_2, \ldots, x_t]
$$
viewed as a chain complex sitting in degrees $(t - 1), \ldots, 0$.
The complex $B_\bullet[1/a_1]$ is isomorphic to the Koszul complex
on $x_2 - a_2/a_1, \ldots, x_t - a_t/a_1$ which is a regular sequence
in $A[1/a_1][x_2, \ldots, x_t]$. Since regular sequences are
Koszul regular, we conclude that the augmentation
$$
\epsilon : B_\bullet \longrightarrow B
$$
is a quasi-isomorphism after inverting $a_1$. Since the homology modules
of the cone $C_\bullet$ on $\epsilon$ are finite $A[x_2, \ldots, x_n]$-modules
and since $C_\bullet$ is bounded,
we conclude that there exists a $c \geq 0$ such that $a_1^c$
annihilates all of these. By
Derived Categories, Lemma \ref{derived-lemma-trick-vanishing-composition}
this implies that, after possibly replacing $c$ by a larger integer,
that $a_1^c$ is zero on $C_\bullet$ in $D(A)$.
The proof is finished once the reader contemplates
the distinguished triangle
$$
B_\bullet \otimes_A^\mathbf{L} M \to
B \otimes_A^\mathbf{L} M \to
C_\bullet \otimes_A^\mathbf{L} M
$$
Namely, the first term is represented by $B_\bullet \otimes_A M$ which
is sitting in homological degrees $(t - 1), \ldots, 0$
in view of the fact that the terms in the Koszul complex $B_\bullet$
are free (and hence flat) $A$-modules. Whence
$\text{Tor}_i^A(B, M) = H_i(C_\bullet \otimes_A^\mathbf{L} M)$
for $i > t - 1$ and this is annihilated by $a_1^c$.
Since $a_1^cB = I^cB$ and since the tor module is a module
over $B$ we conclude.
\end{proof}

\end{document}
